% ============================================================
%  Seção 2.3 — Processamento de Linguagem Natural (PLN)
%  Capítulo 2 — Fundamentação Teórica
% ============================================================

\section{Processamento de Linguagem Natural}
\label{sec:pln}

O Processamento de Linguagem Natural (PLN) é o campo de pesquisa que investiga e propõe métodos e sistemas para o tratamento computacional da linguagem humana \cite{caseli2024cap1}. Situado na interseção entre a Ciência da Computação, a Inteligência Artificial e a Linguística, o PLN abrange desde tarefas de baixo nível --- como segmentação de texto em palavras e identificação de classes gramaticais --- até tarefas de alta complexidade semântica, como tradução automática, resposta a perguntas e análise de sentimento \cite{jurafsky2026}. Para os fins deste trabalho, o PLN constitui o conjunto de técnicas que viabiliza a extração automatizada de sentimento a partir de publicações textuais do X/Twitter relacionadas ao mercado financeiro brasileiro, transformando dados não estruturados em classificações de polaridade que expressam o humor coletivo dos investidores.

% ============================================================
%  Subseção 2.3.1 — Conceitos Fundamentais e Paradigmas
% ============================================================

\subsection{Conceitos Fundamentais e Paradigmas}
\label{subsec:pln_paradigmas}

A linguagem natural apresenta desafios computacionais que a distinguem fundamentalmente de outros tipos de dados: é ambígua, dependente de contexto, culturalmente situada e semanticamente densa \cite{caseli2024cap1}. Uma mesma sequência de palavras pode admitir interpretações radicalmente distintas conforme o contexto discursivo em que é produzida, e o significado de uma expressão frequentemente depende de conhecimento de mundo que não está explicitamente codificado no texto. Essas propriedades fazem com que o processamento completo e adequado da língua continue sendo um dos problemas abertos mais desafiadores da Inteligência Artificial \cite{caseli2024cap1,jurafsky2026}.

Ao longo de sua história, o PLN organizou-se em torno de três grandes paradigmas \cite{caseli2024cap1}. O \textbf{paradigma simbólico}, dominante até a década de 1980, representa o conhecimento linguístico de forma explícita por meio de léxicos, gramáticas e regras lógicas construídas manualmente por especialistas. Embora ofereça alta interpretabilidade e precisão em domínios bem definidos, essa abordagem não escala para a diversidade e a variabilidade da linguagem em uso real. O \textbf{paradigma estatístico}, que ganhou proeminência nos anos 1990 com o aumento da capacidade computacional e a disponibilidade de grandes corpora digitais, substituiu as regras explícitas por modelos probabilísticos aprendidos a partir de dados. Fenômenos linguísticos passaram a ser representados em termos de frequências de coocorrência e distribuições condicionais, com destaque para os modelos de \textit{n-gramas} e os classificadores baseados em aprendizado de máquina raso (\textit{shallow learning}) \cite{jurafsky2026}. O \textbf{paradigma neural}, que se consolidou a partir de meados dos anos 2010, introduziu arquiteturas de redes neurais profundas capazes de aprender representações distribuídas da língua diretamente a partir de texto bruto, sem a necessidade de engenharia de \textit{features} manual. Esse paradigma atingiu seu apogeu com o desenvolvimento dos modelos de linguagem baseados em \textit{transformers} \cite{jurafsky2026}, que estabeleceram novos estados da arte em praticamente todas as tarefas de PLN e são hoje a base dos sistemas de maior desempenho em análise de sentimento financeiro \cite{santos2023finbert}.

Na prática, aplicações contemporâneas de PLN frequentemente adotam abordagens \textbf{híbridas}, combinando o conhecimento explícito do paradigma simbólico --- como léxicos especializados de domínio --- com a capacidade generalizante dos modelos neurais \cite{caseli2024cap1,kearney2014}. No domínio financeiro em português, essa hibridização é especialmente relevante dado o vocabulário técnico específico do mercado brasileiro, que pode não estar adequadamente representado em modelos de linguagem treinados em corpora genéricos \cite{sousa2025,santos2023finbert}.

% ============================================================
%  Subseção 2.3.2 — Pipeline de Processamento Textual
% ============================================================

\subsection{Pipeline de Processamento Textual}
\label{subsec:pln_pipeline}

A extração de informação semântica a partir de textos brutos segue, em geral, um conjunto de etapas encadeadas conhecido como \textit{pipeline} de PLN. Cada etapa transforma a representação do texto, progressivamente reduzindo o ruído e aumentando a estrutura disponível para as etapas subsequentes \cite{jurafsky2026,caseli2024cap1}. No contexto de análise de sentimento de publicações do X/Twitter, as principais etapas são as seguintes.

A \textbf{coleta e armazenamento} constitui o ponto de entrada do \textit{pipeline}. No caso de dados do X/Twitter, isso envolve o acesso à API da plataforma e a estruturação dos dados brutos --- que chegam em formato JSON semiestruturado --- em um formato adequado para processamento \cite{freitas2024cap13}. Metadados como data de publicação, identificador do usuário e métricas de engajamento (\textit{likes}, \textit{retweets}) são preservados para análises subsequentes.

A \textbf{tokenização} é o processo de segmentar o texto em unidades mínimas de processamento, denominadas \textit{tokens} \cite{jurafsky2026}. Em textos de redes sociais, a tokenização apresenta desafios particulares: \textit{hashtags}, \textit{cashtags} (e.g., \texttt{\$PETR4}), menções (\texttt{@usuário}), emojis e abreviações informais são elementos que tokenizadores projetados para texto formal tendem a tratar de maneira inadequada. Tokenizadores especializados para o domínio de redes sociais, ou os próprios tokenizadores internos dos modelos \textit{transformer} (baseados em \textit{subword tokenization} como \textit{WordPiece} ou \textit{Byte-Pair Encoding}), são mais adequados para esse tipo de dado \cite{jurafsky2026,santos2023finbert}.

A \textbf{normalização e limpeza} compreende um conjunto de transformações que reduzem a variabilidade superficial do texto sem alterar seu conteúdo semântico relevante. Para textos financeiros do Twitter, isso tipicamente inclui: remoção de URLs, conversão para minúsculas, tratamento de caracteres especiais e normalização de entidades financeiras (\textit{tickers}, siglas de índices). Em abordagens baseadas em léxicos, etapas adicionais como \textbf{remoção de palavras funcionais} (\textit{stopwords}), \textbf{lematização} --- redução de formas flexionadas à forma canônica da palavra --- e \textbf{stemização} --- truncamento do sufixo para obtenção do radical --- são aplicadas para reduzir a dimensionalidade do vocabulário \cite{caseli2024cap1,sousa2025}. Em modelos \textit{transformer} como o BERT, essas etapas são parcialmente dispensáveis, pois o modelo aprende representações contextualizadas que capturam variações morfológicas automaticamente \cite{jurafsky2026}.

A \textbf{representação vetorial} (\textit{embedding}) é a etapa em que o texto processado é convertido em representações numéricas que capturam relações semânticas entre palavras e sentenças \cite{jurafsky2026}. A evolução nessa etapa reflete a progressão dos paradigmas do PLN: das representações esparsas baseadas em frequências (\textit{bag-of-words}, TF-IDF) às representações densas aprendidas por modelos neurais (\textit{word2vec}, GloVe) e, mais recentemente, às representações contextualizadas produzidas por modelos de linguagem de grande escala baseados em \textit{transformers}, em que o vetor de uma palavra é função de todo o contexto em que ela ocorre \cite{jurafsky2026,santos2023finbert}.

A \textbf{classificação ou extração} constitui a etapa final do \textit{pipeline}, em que as representações vetoriais são utilizadas como entrada para modelos que realizam a tarefa de interesse --- no caso deste trabalho, a classificação do sentimento de cada publicação em polaridades (positiva, negativa ou neutra) \cite{claro2024cap22}. A qualidade dessa etapa depende criticamente da adequação das etapas anteriores e da representatividade dos dados de treinamento do modelo em relação ao domínio e à língua-alvo \cite{kearney2014,loughran2011}.

% ============================================================
%  Subseção 2.3.3 — Modelos de Linguagem e Transformers
% ============================================================

\subsection{Modelos de Linguagem e Transformers}
\label{subsec:transformers}

A arquitetura \textit{Transformer}, introduzida por \citeauthor{vaswani2017} (\citeyear{vaswani2017}) no artigo ``Attention Is All You Need'', representou uma ruptura fundamental no campo do PLN. Diferentemente das arquiteturas recorrentes (RNN, LSTM) que processavam sequências de forma serial, o \textit{Transformer} opera em paralelo sobre toda a sequência de entrada por meio de um mecanismo de \textbf{atenção} (\textit{self-attention}), que pondera a relevância de cada \textit{token} em relação a todos os outros \cite{jurafsky2026}. Isso permitiu tanto um treinamento significativamente mais eficiente em grandes corpora quanto uma captura mais rica de dependências de longo alcance na língua.

O modelo BERT (\textit{Bidirectional Encoder Representations from Transformers}), proposto por \citeauthor{devlin2019} (\citeyear{devlin2019}), consolidou o paradigma de \textbf{pré-treinamento e ajuste fino} (\textit{pre-training and fine-tuning}) como o estado da arte em PLN. Na fase de pré-treinamento, o BERT é treinado em tarefas auto-supervisionadas sobre grandes volumes de texto não rotulado --- \textit{Masked Language Modeling} (MLM), em que o modelo aprende a prever tokens mascarados com base no contexto bidirecional, e \textit{Next Sentence Prediction} (NSP) --- adquirindo representações linguísticas ricas e generalizáveis \cite{jurafsky2026}. Na fase de ajuste fino (\textit{fine-tuning}), o modelo pré-treinado é especializado para uma tarefa específica --- como classificação de sentimento --- com um número relativamente pequeno de exemplos rotulados, transferindo o conhecimento linguístico adquirido na fase anterior \cite{santos2023finbert}.

Para o domínio financeiro em inglês, \citeauthor{araci2019} (\citeyear{araci2019}) desenvolveram o FinBERT, um modelo BERT especializado por meio de ajuste fino em textos financeiros, que demonstrou desempenho superior ao BERT genérico em tarefas de análise de sentimento financeiro. Essa abordagem evidenciou que a especialização de domínio --- o pré-treinamento ou ajuste fino com textos específicos da área --- é um fator determinante para o desempenho em linguagem técnica \cite{araci2019,santos2023finbert}.

No contexto brasileiro, o modelo \textbf{BERTimbau} \cite{souza2020} forneceu a base para a especialização em português: trata-se de um BERT pré-treinado sobre um corpus de 2,68 bilhões de palavras em português, cobrindo fontes jornalísticas, enciclopédicas e da \textit{web}. A partir do BERTimbau, \citeauthor{santos2023finbert} (\citeyear{santos2023finbert}) desenvolveram o \textbf{FinBERT-PT-BR}, realizando ajuste fino sobre 1,4 milhão de sentenças de notícias financeiras em português coletadas do Valor Econômico, Exame e InfoMoney entre 2006 e 2022. O modelo resultante --- e sua variante para classificação de sentimento, o \textbf{SentFinBERT-PT-BR} --- apresentou desempenho superior aos modelos anteriores do estado da arte em português para a tarefa de análise de sentimento de textos financeiros, sendo adotado neste trabalho como a principal ferramenta de extração de polaridade das publicações do X/Twitter relacionadas aos índices da B3.

